\documentclass[10pt,a4paper]{article}

% --- Paquets essentiels ---
\usepackage[utf8]{inputenc}
\usepackage[T1]{fontenc}
\usepackage[french]{babel}
\usepackage{lmodern}
\usepackage{geometry}
\geometry{top=1.2cm, bottom=1.3cm, left=1.3cm, right=1.3cm}
\usepackage[table]{xcolor}
\usepackage{booktabs}
\usepackage{tabularx}
\usepackage{array}
\usepackage{enumitem}
\usepackage{hyperref}
\usepackage{fancyhdr}
\usepackage{titlesec}
\usepackage{tikz}
\usetikzlibrary{shapes,arrows.meta,positioning,fit,backgrounds,calc}

% --- Couleurs ---
\definecolor{primary}{RGB}{24, 43, 73}       % Bleu marine
\definecolor{secondary}{RGB}{185, 120, 30}   % Doré
\definecolor{syncblue}{RGB}{20, 90, 160}     % Bleu synchrone
\definecolor{asyncorange}{RGB}{200, 80, 10}  % Orange asynchrone
\definecolor{rowgray}{RGB}{246, 248, 252}

% --- Configuration des titres ---
\titleformat{\section}{\color{primary}\large\bfseries}{\thesection.}{0.5em}{}[\titlerule]
\titlespacing*{\section}{0pt}{7pt}{4pt}
\titleformat{\subsection}{\color{secondary}\normalsize\bfseries}{\thesubsection.}{0.4em}{}
\titlespacing*{\subsection}{0pt}{4pt}{2pt}

% --- En-têtes et pieds de page ---
\pagestyle{fancy}
\fancyhf{}
\fancyhead[L]{\footnotesize\color{gray}\textbf{Golden House} -- Architecture Microservices}
\fancyhead[R]{\footnotesize\color{gray}Polytech Montpellier -- IWA 2}
\fancyfoot[C]{\footnotesize\thepage\ / 4}
\renewcommand{\headrulewidth}{0.4pt}
\renewcommand{\footrulewidth}{0.4pt}

\hypersetup{colorlinks=true, linkcolor=primary, urlcolor=syncblue}

\begin{document}

% =============================================================================
% EN-TÊTE COMPACT (Page 1)
% =============================================================================
\noindent
\begin{minipage}[c]{0.65\textwidth}
    {\LARGE\bfseries\color{primary} GOLDEN HOUSE}\\[2pt]
    {\large\bfseries\color{secondary} Architecture des microservices}\\[3pt]
    {\footnotesize\textbf{Auteurs :} Justin Chapon, Nabil Bennacer, Tom Cantillon, Salma El Annaby, Manuela Zapata Quirós}
\end{minipage}
\hfill
\begin{minipage}[c]{0.32\textwidth}
    \raggedleft\footnotesize
    \textbf{Polytech Montpellier} -- IG5\\
    Module : \textbf{IWA 2}\\
    Semestre 9 -- Septembre 2026
\end{minipage}

\vspace{3pt}
\hrule height 1.2pt
\vspace{5pt}

% =============================================================================
% SECTION 1 : SCHÉMA D'ARCHITECTURE (Page 1)
% =============================================================================
\section{Schéma des microservices}

\begin{center}
\begin{tikzpicture}[
    scale=0.88, every node/.style={transform shape},
    box/.style={rectangle, draw=black!70, rounded corners=3pt, text width=2.6cm, align=center, inner sep=4pt, font=\sffamily\scriptsize},
    client/.style={box, fill=gray!15, text width=2.9cm},
    gw/.style={box, fill=purple!12, draw=purple!80!black, text width=5.5cm},
    gwsmall/.style={box, fill=purple!12, draw=purple!80!black, text width=2.8cm},
    service/.style={box, fill=blue!8, draw=blue!70!black},
    billing/.style={box, fill=yellow!22, draw=orange!90!black},
    batch/.style={box, fill=green!10, draw=green!60!black},
    ext/.style={box, fill=red!10, draw=red!70!black, text width=2.4cm},
    db/.style={cylinder, cylinder uses custom fill, cylinder end fill=gray!30, cylinder body fill=gray!15, shape border rotate=90, draw=black!70, aspect=0.25, text width=2.8cm, align=center, font=\sffamily\scriptsize},
    sync/.style={->, >=Stealth, thick, color=syncblue},
    async/.style={->, >=Stealth, thick, dashed, color=asyncorange},
    route/.style={->, >=Stealth, color=gray!70!black}
]

% 1. Clients (Y=11.2)
\node[client] (mobile) at (3.8, 11.2) {\textbf{App Mobile Acheteur}\\\tiny React Native};
\node[client] (backoffice) at (8.2, 11.2) {\textbf{Backoffice Web}\\\tiny React / TypeScript};
\node[client] (chatbot) at (14.8, 11.2) {\textbf{Chatbot IA}\\\tiny Assistant client (LLM)};

% 2. Gateways (Y=9.8)
\node[gw] (apigw) at (6.0, 9.8) {\textbf{API Gateway} (Spring Cloud Gateway)\\\tiny Routage dynamique, Sécurité JWT \& Rate Limiting};
\node[gwsmall] (mcpgw) at (14.8, 9.8) {\textbf{MCP Gateway}\\\tiny Model Context Protocol};

% 3. Row 1 : Tunnel d'achat direct (Y=7.4)
\node[service] (user) at (1.6, 7.4) {\textbf{User Service}\\\tiny Profils, Adresses, Auth};
\node[service] (order) at (6.0, 7.4) {\textbf{Order Service}\\\tiny Cycle de commande};
\node[service] (cart) at (10.4, 7.4) {\textbf{Cart Service}\\\tiny Panier \& Totaux};
\node[service] (catalog) at (14.8, 7.4) {\textbf{Catalog Service}\\\tiny Meubles, Stocks};

% 4. Row 2 : Facturation, Paiement, Notification, Fournisseurs (Y=4.8)
\node[billing] (billing) at (1.6, 4.8) {\textbf{Billing Service}\\\tiny Facturation légale PDF};
\node[service] (payment) at (6.0, 4.8) {\textbf{Payment Service}\\\tiny Sessions Stripe Checkout};
\node[service] (notif) at (10.4, 4.8) {\textbf{Notification Service}\\\tiny Emails (SMTP) \& Push};
\node[service] (supplier) at (14.8, 4.8) {\textbf{Supplier Service}\\\tiny Réapprovisionnements};

% 5. Row 3 : Batches Spark, Tiers Stripe, PostgreSQL, Avis (Y=2.2)
\node[batch] (batches) at (1.6, 2.2) {\textbf{Batches Spark}\\\tiny KPIs ventes \& Reco ALS};
\node[ext] (stripe) at (6.0, 2.2) {\textbf{Stripe API}\\\tiny Prestataire bancaire CB};
\node[db] (postgres) at (10.4, 2.2) {\textbf{PostgreSQL}\\\tiny Schémas dédiés ACID};
\node[service] (review) at (14.8, 2.2) {\textbf{Review Service}\\\tiny Avis \& Notes produits};

% Routes Gateways
\draw[route] (mobile) -- (apigw);
\draw[route] (backoffice) -- (apigw);
\draw[route] (chatbot) -- (mcpgw);

\draw[route] (apigw.south) -- (user.north);
\draw[route] (apigw.south) -- (order.north);
\draw[route] (apigw.south) -- (cart.north);
\draw[route] (apigw.south) -- (catalog.north);
\draw[route] (mcpgw.south) -- (order.north);
\draw[route] (mcpgw.south) -- (catalog.north);

% Liens Synchrones (Bleu continu)
\draw[sync] (order.west) -- node[above, font=\tiny\bfseries] {Valide adr.} (user.east);
\draw[sync] (order.east) -- node[above, font=\tiny\bfseries] {Détail panier} (cart.west);
\draw[sync] (cart.east) -- node[above, font=\tiny\bfseries] {Prix / Stock} (catalog.west);
\draw[sync] (order.north) to[out=25, in=155] node[above, font=\tiny\bfseries] {Réserve stock} (catalog.north);
\draw[sync] (order.south) -- node[right, font=\tiny\bfseries] {Init session} (payment.north);
\draw[sync] (payment.south) -- node[right, font=\tiny\bfseries] {Checkout API} (stripe.north);
\draw[sync] (review.east) to[out=30, in=330] node[right, font=\tiny\bfseries] {Vérif produit} (catalog.east);

% Liens Asynchrones (Orange tireté)
\draw[async] (stripe.north west) to[out=135, in=225] node[left, font=\tiny\bfseries] {Webhook} (payment.south west);
\draw[async] (payment.north west) to[out=120, in=240] node[left, font=\tiny\bfseries] {Paiement OK} (order.south west);
\draw[async] (payment.west) -- node[above, font=\tiny\bfseries] {OrderPaid} (billing.east);
\draw[async] (billing.south east) to[out=300, in=240] node[below, font=\tiny\bfseries] {Facture PDF (\texttt{InvoiceGenerated})} (notif.south west);
\draw[async] (payment.east) -- node[above, font=\tiny\bfseries] {Achat OK} (notif.west);
\draw[async] (order.south east) -- node[above, sloped, font=\tiny\bfseries] {Statut cmd} (notif.north west);
\draw[async] (catalog.south west) -- node[above, sloped, font=\tiny\bfseries] {Seuil stock} (notif.north east);
\draw[async] (catalog.south) -- node[right, font=\tiny\bfseries] {Réappro} (supplier.north);
\draw[async] (supplier.north east) to[out=45, in=315] node[right, font=\tiny\bfseries] {Stock reçu} (catalog.south east);
\draw[async] (review.north) -- node[right, font=\tiny\bfseries] {Note moy.} (supplier.south);

% Base de données (Dotted gray)
\draw[->, >=Stealth, dotted, thick, gray!70] (batches.south) to[out=315, in=190] (postgres.west);
\draw[->, >=Stealth, dotted, thick, gray!70] (notif.south) -- (postgres.north);
\draw[->, >=Stealth, dotted, thick, gray!70] (supplier.south west) -- (postgres.north east);
\draw[->, >=Stealth, dotted, thick, gray!70] (review.west) -- (postgres.east);

% Légende
\node[draw=black!50, fill=white, rounded corners=2pt, inner sep=3pt, font=\tiny, align=center] at (8.2, 0.5) {
    \textbf{Légende :} \quad
    \textcolor{syncblue}{\textbf{------}} \textbf{Synchrone} (REST / HTTP bloquant) \qquad
    \textcolor{asyncorange}{\textbf{- - - -}} \textbf{Asynchrone} (Events / Webhooks découplés) \qquad
    \textcolor{gray!70!black}{\textbf{$\longrightarrow$}} Routage Gateway \qquad
    \textcolor{gray!70}{\textbf{$\cdots\cdots$}} Persistance PostgreSQL
};

\end{tikzpicture}
\end{center}

\vspace{-4pt}

% =============================================================================
% SECTION 2 : RÈGLE DE DÉCISION (Page 1)
% =============================================================================
\section{Synchrone et Asynchrone}

\begin{itemize}[leftmargin=*, nosep]
    \item \textbf{\color{syncblue}Synchrone (REST / HTTP)} : Requêtes bloquantes exigeant une réponse immédiate pour finaliser l'action courante (lecture catalogue, calcul panier, validation adresse, réservation stock, session Stripe).
    \item \textbf{\color{asyncorange}Asynchrone (Events / Webhooks)} : Traitements découplés dans le temps et tolérants aux pannes sans bloquer l'acheteur (webhook Stripe, confirmation paiement, génération facture PDF, notifications email/push, réassort fournisseur, recalcul note).
\end{itemize}

\newpage

% =============================================================================
% SECTION 3 : MATRICE DES COMMUNICATIONS (Page 2)
% =============================================================================
\section{Communications entre services}

\begin{center}
\renewcommand{\arraystretch}{1.18}
\small
\rowcolors{2}{rowgray}{white}
\begin{tabularx}{\textwidth}{@{} l l c p{4.0cm} X @{}}
\toprule[1.2pt]
\textbf{Source} & \textbf{Cible} & \textbf{Type} & \textbf{Protocole / Événement} & \textbf{Description} \\
\midrule[0.8pt]
\textbf{Clients (Mobile/Web)} & API Gateway & \textcolor{syncblue}{\textbf{Sync}} & HTTP/2 / REST & Point d'entrée unique, terminaison TLS, validation JWT. \\ \hline
\textbf{API Gateway} & Microservices & \textcolor{syncblue}{\textbf{Sync}} & HTTP REST interne & Routage inverse vers les services cibles selon l'URI. \\ \hline
\textbf{Chatbot IA} & MCP Gateway & \textcolor{syncblue}{\textbf{Sync}} & JSON-RPC / REST & Traduction du langage naturel en requêtes microservices. \\ \hline
\textbf{Order} & User & \textcolor{syncblue}{\textbf{Sync}} & \texttt{GET /api/users/\{id\}} & Contrôle bloquant de l'existence client et adresse. \\ \hline
\textbf{Order} & Cart & \textcolor{syncblue}{\textbf{Sync}} & \texttt{GET /api/cart/\{id\}} & Récupération synchrone du snapshot panier à figer. \\ \hline
\textbf{Cart} & Catalog & \textcolor{syncblue}{\textbf{Sync}} & \texttt{GET /api/products/\{id\}} & Contrôle du stock disponible et du prix unitaire. \\ \hline
\textbf{Order} & Catalog & \textcolor{syncblue}{\textbf{Sync}} & \texttt{POST /products/reserve} & Verrouillage atomique des stocks avant confirmation. \\ \hline
\textbf{Order} & Payment & \textcolor{syncblue}{\textbf{Sync}} & \texttt{POST /payments/init} & Création de la session d'encaissement et lien Stripe. \\ \hline
\textbf{Payment} & Stripe & \textcolor{syncblue}{\textbf{Sync}} & HTTPS / REST & Appel synchrone à Stripe pour initialiser le Checkout. \\ \hline
\textbf{Review} & Catalog & \textcolor{syncblue}{\textbf{Sync}} & \texttt{GET /api/products/\{id\}} & Vérification de l'existence du meuble noté. \\ \hline
\textbf{Stripe} & Payment & \textcolor{asyncorange}{\textbf{Async}} & Webhook \texttt{checkout.session} & Signal bancaire externe reçu hors session utilisateur. \\ \hline
\textbf{Payment} & Order & \textcolor{asyncorange}{\textbf{Async}} & Event \texttt{PaymentConfirmed} & Notification de règlement : commande passe à \texttt{PAID}. \\ \hline
\textbf{Payment} & Billing & \textcolor{asyncorange}{\textbf{Async}} & Event \texttt{OrderPaidEvent} & Déclenchement de la facture légale sans bloquer l'achat. \\ \hline
\textbf{Billing} & Notification & \textcolor{asyncorange}{\textbf{Async}} & Event \texttt{InvoiceGenerated} & Envoi de l'URL du PDF de facture pour transmission client. \\ \hline
\textbf{Order} & Notification & \textcolor{asyncorange}{\textbf{Async}} & Event \texttt{OrderStatusChanged} & Notification du statut client (expédiée, livrée). \\ \hline
\textbf{Payment} & Notification & \textcolor{asyncorange}{\textbf{Async}} & Event \texttt{PaymentSucceeded} & Envoi immédiat de l'email de confirmation et reçu. \\ \hline
\textbf{Catalog} & Notification & \textcolor{asyncorange}{\textbf{Async}} & Event \texttt{StockBelowThreshold} & Alerte gestionnaire si stock sous le seuil critique. \\ \hline
\textbf{Catalog} & Supplier & \textcolor{asyncorange}{\textbf{Async}} & Event \texttt{RestockNeededEvent} & Demande automatique de commande de réapprovisionnement. \\ \hline
\textbf{Supplier} & Catalog & \textcolor{asyncorange}{\textbf{Async}} & Event \texttt{StockReceivedEvent} & Augmentation des stocks à la réception des meubles. \\ \hline
\textbf{Review} & Supplier & \textcolor{asyncorange}{\textbf{Async}} & Event \texttt{RatingAggregated} & Recalcul de la note moyenne du fabricant associé. \\ \hline
\textbf{Batches} & Order / Catalog & \textcolor{asyncorange}{\textbf{Async}} & Jobs Spark SQL & Traitements nocturnes pour KPIs et modèles ALS. \\
\bottomrule[1.2pt]
\end{tabularx}
\end{center}

\newpage

% =============================================================================
% SECTION 4 : POINTS D'ENTRÉE & MODÈLES DE DONNÉES (Pages 3 et 4)
% =============================================================================
\section{Points d'entrée et données (1/2)}

\subsection{Gateways}
\begin{itemize}[leftmargin=*, nosep]
    \item \textbf{Points d'entrée} :
    \begin{itemize}[nosep]
        \item \texttt{ANY /api/v1/**} : Routage dynamique vers microservices, terminaison TLS, validation token JWT.
        \item \texttt{POST /mcp/v1/query} : Requête contextuelle Chatbot IA (Model Context Protocol).
    \end{itemize}
    \item \textbf{Données} : \textit{Stateless}. Cache Redis pour la révocation des tokens JWT et le rate limiting.
\end{itemize}

\vspace{2pt}
\subsection{User Service}
\begin{itemize}[leftmargin=*, nosep]
    \item \textbf{Points d'entrée} :
    \begin{itemize}[nosep]
        \item \texttt{POST /api/users/register} : Inscription nouvel utilisateur.
        \item \texttt{POST /api/users/login} : Authentification et génération du token JWT.
        \item \texttt{GET | PUT /api/users/\{id\}} : Consultation et mise à jour du profil utilisateur.
        \item \texttt{GET | POST /api/users/\{id\}/addresses} : Gestion des adresses de facturation et de livraison.
    \end{itemize}
    \item \textbf{Modèle PostgreSQL} :
    \begin{itemize}[nosep]
        \item \texttt{\textbf{users}} : \texttt{id} (UUID, PK), \texttt{email} (VARCHAR, UK), \texttt{username} (VARCHAR, UK), \texttt{password\_hash}, \texttt{first\_name}, \texttt{last\_name}, \texttt{phone}, \texttt{role} (role\_enum), \texttt{currency} (currency\_enum), \texttt{stripe\_customer\_id} (UK), \texttt{is\_active} (BOOL).
        \item \texttt{\textbf{addresses}} : \texttt{id} (UUID, PK), \texttt{user\_id} (UUID, FK users), \texttt{type} (address\_type), \texttt{street}, \texttt{postal\_code}, \texttt{city}, \texttt{country}, \texttt{is\_default} (BOOL).
    \end{itemize}
\end{itemize}

\vspace{2pt}
\subsection{Catalog Service}
\begin{itemize}[leftmargin=*, nosep]
    \item \textbf{Points d'entrée} :
    \begin{itemize}[nosep]
        \item \texttt{GET /api/products} : Liste paginée des meubles avec filtres (catégorie, matière, couleur, prix).
        \item \texttt{GET /api/products/\{id\}} : Fiche détaillée d'un meuble (dimensions, stock disponible).
        \item \texttt{POST /api/products} : Création / mise à jour produit (réservé Backoffice).
        \item \texttt{POST /api/products/reserve} : Réservation atomique de quantité lors de la commande.
    \end{itemize}
    \item \textbf{Modèle PostgreSQL} :
    \begin{itemize}[nosep]
        \item \texttt{\textbf{products}} : \texttt{id} (UUID, PK), \texttt{nom} (VARCHAR), \texttt{prix} (NUMERIC), \texttt{categorie} (category\_enum), \texttt{type\_matiere} (material\_enum), \texttt{couleur} (VARCHAR), \texttt{largeur/hauteur/prof.} (INT), \texttt{stock} (INT), \texttt{seuil\_alerte} (INT), \texttt{images} (TEXT[]), \texttt{fournisseur\_id} (UUID, FK), \texttt{est\_actif} (BOOL).
    \end{itemize}
\end{itemize}

\vspace{2pt}
\subsection{Cart Service}
\begin{itemize}[leftmargin=*, nosep]
    \item \textbf{Points d'entrée} :
    \begin{itemize}[nosep]
        \item \texttt{GET /api/cart/\{userId\}} : Récupération du panier actif et calcul des sous-totaux.
        \item \texttt{POST /api/cart/items} : Ajout d'un article dans le panier.
        \item \texttt{PUT /api/cart/items/\{id\}} : Modification de quantité d'un article.
        \item \texttt{DELETE /api/cart/items/\{id\}} : Suppression d'un article ; \texttt{DELETE /api/cart/\{id\}} : Vidage du panier.
    \end{itemize}
    \item \textbf{Modèle PostgreSQL} :
    \begin{itemize}[nosep]
        \item \texttt{\textbf{carts}} : \texttt{id} (UUID, PK), \texttt{user\_id} (UUID, FK-UQ users), \texttt{code\_promo} (VARCHAR).
        \item \texttt{\textbf{cart\_items}} : \texttt{id} (UUID, PK), \texttt{cart\_id} (UUID, FK carts), \texttt{product\_id} (UUID, FK products), \texttt{quantite} (INT).
    \end{itemize}
\end{itemize}

\vspace{2pt}
\subsection{Order Service}
\begin{itemize}[leftmargin=*, nosep]
    \item \textbf{Points d'entrée} :
    \begin{itemize}[nosep]
        \item \texttt{POST /api/orders} : Création de commande depuis le panier, validation adresse et réservation stock.
        \item \texttt{GET /api/orders/\{id\}} : Détail d'une commande (lignes, montants, adresses, statuts).
        \item \texttt{GET /api/orders/user/\{userId\}} : Historique complet des commandes de l'acheteur.
        \item \texttt{PUT /api/orders/\{id\}/status} : Mise à jour du statut (\texttt{PENDING}, \texttt{PAID}, \texttt{SHIPPED}, \texttt{DELIVERED}).
    \end{itemize}
    \item \textbf{Modèle PostgreSQL} :
    \begin{itemize}[nosep]
        \item \texttt{\textbf{orders}} : \texttt{id} (UUID, PK), \texttt{order\_number} (VARCHAR, UK), \texttt{user\_id} (UUID, FK), \texttt{statut} (order\_status\_enum), \texttt{statut\_paiement} (payment\_enum), \texttt{stripe\_payment\_intent\_id} (VARCHAR, UK), \texttt{shipping\_address\_id} (UUID, FK), \texttt{billing\_address\_id} (UUID, FK), \texttt{montant\_total} (NUMERIC), \texttt{devise} (currency\_enum).
        \item \texttt{\textbf{order\_items}} : \texttt{id} (UUID, PK), \texttt{order\_id} (UUID, FK orders), \texttt{product\_id} (UUID, FK products), \texttt{nom\_produit} (VARCHAR), \texttt{quantite} (INT), \texttt{prix\_unitaire} (NUMERIC), \texttt{total\_ligne} (NUMERIC).
    \end{itemize}
\end{itemize}

\newpage

% =============================================================================
% SECTION 5 : POINTS D'ENTRÉE & MODÈLES DE DONNÉES (Page 4)
% =============================================================================
\section{Points d'entrée et données (2/2)}

\subsection{Payment Service}
\begin{itemize}[leftmargin=*, nosep]
    \item \textbf{Points d'entrée} :
    \begin{itemize}[nosep]
        \item \texttt{POST /api/payments/init} : Création de la session Stripe Checkout et renvoi de l'URL de paiement.
        \item \texttt{POST /api/payments/webhook} : Réception asynchrone des événements de paiement Stripe.
        \item \texttt{GET /api/payments/order/\{orderId\}} : Consultation de l'état de la transaction financière.
    \end{itemize}
    \item \textbf{Modèle PostgreSQL} :
    \begin{itemize}[nosep]
        \item \texttt{\textbf{payments}} : \texttt{id} (UUID, PK), \texttt{order\_id} (UUID, UK), \texttt{stripe\_session\_id}, \texttt{stripe\_payment\_intent\_id}, \texttt{amount} (NUMERIC), \texttt{currency}, \texttt{status} (payment\_enum), \texttt{paid\_at} (TIMESTAMPTZ).
    \end{itemize}
\end{itemize}

\vspace{2pt}
\subsection{Billing Service}
\begin{itemize}[leftmargin=*, nosep]
    \item \textbf{Rôle} : Consommateur asynchrone de l'événement \texttt{OrderPaidEvent}. Attribue un numéro fiscal chronologique unique, calcule la ventilation HT / TVA / TTC, et compile le document légal PDF stocké sur S3 / MinIO.
    \item \textbf{Points d'entrée} :
    \begin{itemize}[nosep]
        \item \texttt{GET /api/invoices/order/\{orderId\}} : Récupération des métadonnées et URL du PDF de facture.
        \item \texttt{GET /api/invoices/\{id\}/pdf} : Téléchargement direct du document PDF certifié.
        \item \texttt{GET /api/invoices/user/\{userId\}} : Liste des factures de l'acheteur pour consultation comptable.
    \end{itemize}
    \item \textbf{Modèle PostgreSQL} :
    \begin{itemize}[nosep]
        \item \texttt{\textbf{invoices}} : \texttt{id} (UUID, PK), \texttt{numero\_facture} (VARCHAR, UK, ex: \texttt{FACT-2026-0001}), \texttt{order\_id} (UUID, FK-UQ orders), \texttt{user\_id} (UUID, FK users), \texttt{montant\_ht} (NUMERIC), \texttt{taux\_tva} (NUMERIC), \texttt{montant\_tva} (NUMERIC), \texttt{montant\_ttc} (NUMERIC), \texttt{url\_pdf} (VARCHAR), \texttt{date\_facturation} (TIMESTAMPTZ).
    \end{itemize}
\end{itemize}

\vspace{2pt}
\subsection{Notification Service}
\begin{itemize}[leftmargin=*, nosep]
    \item \textbf{Points d'entrée} :
    \begin{itemize}[nosep]
        \item \texttt{POST /api/notifications/send} : Envoi direct déclenché par les services internes.
        \item \texttt{GET /api/notifications/user/\{userId\}} : Centre de notifications in-app pour l'application mobile.
        \item \texttt{PUT /api/notifications/\{id\}/read} : Marquage de lecture d'une notification.
    \end{itemize}
    \item \textbf{Modèle PostgreSQL} :
    \begin{itemize}[nosep]
        \item \texttt{\textbf{notifications}} : \texttt{id} (UUID, PK), \texttt{user\_id} (UUID, FK users), \texttt{titre} (VARCHAR), \texttt{type} (notif\_type\_enum : EMAIL, PUSH), \texttt{est\_lu} (BOOL).
    \end{itemize}
\end{itemize}

\vspace{2pt}
\subsection{Supplier Service}
\begin{itemize}[leftmargin=*, nosep]
    \item \textbf{Points d'entrée} :
    \begin{itemize}[nosep]
        \item \texttt{GET /api/suppliers} : Annuaire des fabricants partenaires et coordonnées.
        \item \texttt{POST /api/suppliers/orders} : Émission d'un ordre d'achat automatique ou manuel suite à alerte stock.
        \item \texttt{PUT /api/suppliers/orders/\{id\}/receive} : Réception de livraison et déclenchement d'incrément de stock.
    \end{itemize}
    \item \textbf{Modèle PostgreSQL} :
    \begin{itemize}[nosep]
        \item \texttt{\textbf{fournisseurs}} : \texttt{id} (UUID, PK), \texttt{nom} (VARCHAR), \texttt{email} (VARCHAR), \texttt{telephone} (VARCHAR), \texttt{adresse} (TEXT).
        \item \texttt{\textbf{supply\_orders}} : \texttt{id} (UUID, PK), \texttt{fournisseur\_id} (UUID, FK), \texttt{statut} (supply\_status\_enum), \texttt{montant\_total} (NUMERIC), \texttt{date\_emission} (TIMESTAMPTZ), \texttt{date\_reception} (TIMESTAMPTZ).
        \item \texttt{\textbf{supply\_order\_items}} : \texttt{id} (UUID, PK), \texttt{supply\_order\_id} (UUID, FK), \texttt{product\_id} (UUID, FK), \texttt{quantite} (INT), \texttt{prix\_unitaire} (NUMERIC).
    \end{itemize}
\end{itemize}

\vspace{2pt}
\subsection{Review Service et Batches Spark}
\begin{itemize}[leftmargin=*, nosep]
    \item \textbf{Review Service (Avis)} :
    \begin{itemize}[nosep]
        \item \textbf{Points d'entrée} : \texttt{POST /api/reviews} (dépôt avis vérifié), \texttt{GET /api/reviews/product/\{id\}} (avis du produit).
        \item \textbf{Modèle PostgreSQL (\texttt{reviews})} : \texttt{id} (UUID, PK), \texttt{user\_id} (UUID, FK), \texttt{product\_id} (UUID, FK), \texttt{titre} (VARCHAR), \texttt{note} (INT 1-5), \texttt{reponse\_commercant} (TEXT).
    \end{itemize}
    \item \textbf{Batches Spark (Analytique \& Recommandation)} :
    \begin{itemize}[nosep]
        \item \textbf{Traitements} : Traitements Spark SQL distribués de nuit (chiffre d'affaires, panier moyen, volumes par catégorie) et entraînement du modèle de recommandation collaborative ALS (Alternating Least Squares).
        \item \textbf{Modèle PostgreSQL} :
        \begin{itemize}[nosep]
            \item \texttt{\textbf{product\_recommendations}} : \texttt{user\_id}, \texttt{product\_id}, \texttt{score} (FLOAT).
            \item \texttt{\textbf{kpis\_ventes\_journalieres}} : \texttt{date}, \texttt{nb\_commandes}, \texttt{chiffre\_affaires\_ht}.
        \end{itemize}
    \end{itemize}
\end{itemize}

\end{document}
